\label{sec:legal}

\subsection{Section~5's Prophylactic Function}

Section~5 preclearance served a function that Section~2 litigation
cannot replicate: it was a prophylactic rule that deterred discriminatory
redistricting before it was enacted, not a remedy applied after the harm
had occurred. Even jurisdictions that believed their redistricting plans
were VRA-compliant had to justify that belief to the DOJ or a federal
court before implementation. This review process had three effects:

\begin{enumerate}
\item \textbf{Deterrence}: Knowing that plans would be scrutinized
  ex ante discouraged covered jurisdictions from drawing maps that
  would fail the non-retrogression standard.

\item \textbf{Documentation}: The preclearance submission process
  required covered jurisdictions to document their redistricting
  criteria, methods, and demographic analysis. This created a
  contemporaneous record that could be used in subsequent litigation.

\item \textbf{Speed}: Administrative preclearance (the DOJ route)
  was available within 60 days; even the slower judicial preclearance
  route produced decisions before the election cycle. Section~2
  litigation, by contrast, typically spans multiple election cycles.
\end{enumerate}

All three of these functions disappeared with \textit{Shelby County}.
No deterrence exists in former covered states because no pre-enactment
review is required. No documentation requirement compels contemporaneous
disclosure of redistricting criteria and methodology. No fast-track
process exists to resolve disputes before an election cycle begins.

\subsection{Algorithmic Redistricting as Transparency Substitute}

\textsc{bisect} cannot restore Section~5 preclearance---that requires
an act of Congress. But it can provide a transparency substitute that
addresses two of Section~5's three functions: deterrence (partially)
and documentation (fully).

\textbf{Documentation}: The \textsc{bisect} pipeline produces a complete
SHA audit chain for every redistricting plan. Every input (Census tract
data, adjacency graph), every intermediate result (bisection sequence,
seed selection), and every output (district assignment file) is
cryptographically hashed and chained. This produces a plan that is
fully reproducible by any party with access to the public Census data
and the published \textsc{bisect} software. In expert witness proceedings,
this means that the \textsc{bisect} plan can be verified independently
by opposing experts, eliminating disputes about whether the algorithmic
result was manipulated.

\textbf{Deterrence (partial)}: A \textsc{bisect} reference plan available
as a published neutral baseline before state legislatures begin drawing
maps creates a deterrent that is weaker than Section~5 (it is not
legally required) but stronger than nothing. Legislators and mapmakers
who know that a neutral algorithmic plan will be produced and introduced
into evidence must draw further from the neutral baseline to achieve
a desired partisan or demographic result---and the further they draw,
the more visible the deviation from the neutral plan becomes.

\subsection{Three Use Cases in Post-\textit{Shelby} Litigation}

The \textsc{bisect} reference plan serves three distinct roles in
post-\textit{Shelby} Section~2 redistricting litigation.

\textbf{Use Case 1: Plaintiff Reference Plan}. In a Section~2 vote-dilution
case, the plaintiff must show not only that the enacted plan dilutes
minority voting strength but that an alternative plan exists that would
provide adequate minority representation without sacrificing other
redistricting criteria. Under \textit{Allen v.\ Milligan}, the plaintiff's
burden includes demonstrating that a majority-minority district is
``reasonably configured''---that is, compact and otherwise consistent
with legitimate redistricting principles. A \textsc{bisect} plan with
VRA mode enabled (\texttt{--weights-override vra-aligned}) provides
exactly this: a reasonably configured alternative plan that creates the
required majority-minority district while satisfying population balance,
contiguity, and compactness requirements. Because the plan is algorithmic
and reproducible, it is resistant to the opposing argument that it was
``drawn to prove a point'' or constructed after the fact to win litigation.

\textbf{Use Case 2: Defendant's Neutral Map Defense}. States defending
their redistricting plans against Section~2 claims can introduce a
\textsc{bisect} reference plan to show that the enacted map does not
deviate materially from what a neutral process would produce. If the
\textsc{bisect} plan and the enacted plan produce similar BVAP
distributions, the defendant can argue that any deviation from a
``perfect'' minority representation outcome is attributable to geography
rather than mapmaker intent. This defense is most compelling when
the state ran \textsc{bisect} as part of its redistricting process
and adopted the plan with minimal modification---demonstrating that
the neutral algorithm, not political manipulation, determined the
outcome.

\textbf{Use Case 3: Court-Appointed Special Master Baseline}. When
courts order remedial redistricting plans---as the Alabama district court
did in \textit{Allen v.\ Milligan} and the North Carolina Supreme Court
did in \textit{Harper v.\ Hall}---they often appoint special masters
to produce alternative plans. \textsc{bisect} is well-suited to this
role: it produces plans quickly, satisfies the traditional redistricting
criteria that courts require of remedial maps (population equality,
contiguity, compactness), and can be configured to satisfy VRA
requirements through the VRA mode. The full audit chain provides
the transparency that a court-ordered plan requires.

\subsection{The Non-Retrogression Analogy}

Section~5's non-retrogression standard asked: does the new plan reduce
minority electoral opportunity relative to the existing plan? The
\textsc{bisect} reference plan enables an analogous comparison in
Section~2 litigation: does the enacted plan reduce minority electoral
opportunity relative to what a neutral, non-partisan process would produce?

This ``neutral baseline'' comparison does not replicate Section~5's
legal effect---Section~2 does not adopt a non-retrogression standard,
and \textit{Gingles} does not require comparison to a neutral baseline.
But it provides a framework for the expert witness to quantify the
\textit{magnitude} of minority vote dilution. A claim that a map
``dilutes minority voting strength'' is more compelling, and more
legally actionable, when it can be attached to a specific quantity:
``the enacted map produces a BVAP of X\% in the relevant district,
while a neutral algorithmic plan produces Y\%, a gap of Z percentage
points attributable entirely to mapmaker choices rather than demographic
geography.'' This quantification---the gap between the neutral baseline
and the enacted plan---is the algorithmic redistricting contribution
to post-\textit{Shelby} Section~2 litigation.
